\chapter{Hardcoded-Assumption Audit and a CI Guard (Phase 1.7)}
\label{chap:p1-7}

\section{Goal}

Chapter~13 of the pilot report found a real bug the hard way: \texttt{11\_make\_figures.py} had the
sample count (``163'') and per-class sizes typed as string literals into its chart labels, so the
moment the script ran on a different sample size it produced a chart confidently mislabeled with the
wrong numbers. That bug survived twelve days and a passing 61-test suite precisely because every run
until then happened to use the one sample size for which the literals were correct. The scale-up plan
(item 3.7) generalizes the lesson: sweep every analysis and plotting path for literal sample counts,
class sizes, grid dimensions, and the patch-grid arithmetic (the 768$\rightarrow$24 image-token math),
replace each with a value read from data or config, and add an automated check that fails if a known-
magic literal ever reappears. This is the phase that designs the bug \emph{class} out, not just the one
instance Chapter~13 caught.

\section{What Was Done}

\subsection{The patch-grid arithmetic, now derived}

The clearest remaining instance was in \texttt{attribution.py}, which hardcoded the cross-attention
heatmap grid as \texttt{PATCH\_GRID = (24, 24)} with \texttt{N\_PATCH\_TOKENS = 576} --- values correct
only for Florence-2's fixed 768$\times$768 processor input. These are now computed at run time from the
\emph{actual} \texttt{pixel\_values} resolution and the vision tower's documented stride
(\texttt{\_patch\_grid}), and the reshape is protected by an assertion that the encoder sequence is long
enough to hold the global token plus every patch token --- so a layout change fails loudly instead of
silently shifting every heatmap by a cell. The one retained constant, \texttt{VISION\_TOWER\_STRIDE =
32}, is an architectural fact (the DaViT tower's $4\times2\times2\times2$ downsampling), not a
resolution assumption. Verified directly: \texttt{\_patch\_grid} returns $(24,24)$ for a $768\times768$
input --- reproducing the frozen grid exactly --- and $(32,24)$ for a $1024\times768$ input, adapting as
intended.

\subsection{The remaining label literals}

Two output labels still carried hardcoded counts and were switched to data-derived values:
\texttt{08\_efficiency\_test.py}'s plot title ``\ldots(163-sample pilot)'' now interpolates the
observed sample count it already computes, and \texttt{06\_visual\_sparsity.py}'s
``\ldots/ 576 cells'' summary now reads the cell count from the actual heatmaps (recorded as an
\texttt{n\_cells} column) rather than assuming 576. Model IDs were checked and are already sourced from
\texttt{config.MODEL\_ID} everywhere they drive behavior; only cosmetic ``Loading Florence-2\ldots'' log
strings mention the name literally, which affects nothing computed.

\subsection{A CI guard that tells code from documentation}

The automated check is a test, \texttt{test\_no\_hardcoded\_literals.py}, that tokenizes every file
under \texttt{scripts/} and \texttt{src/xai\_pilot/} and fails if a forbidden numeric literal ---
\{163, 3004, 576, 768\}: the pilot/test sample counts and the patch-grid dimensions --- appears as an
actual \texttt{NUMBER} token in code. Using the tokenizer rather than a text grep is the point: comments
and docstrings that legitimately \emph{describe} the pilot's 163 samples or explain why $576 = 24\times
24$ are not \texttt{NUMBER} tokens and are correctly ignored, so the guard catches a real regression (a
literal typed into an expression) without firing on prose that mentions the same number. The check runs
as part of the ordinary \texttt{pytest} suite, so it gates every change the way any other test does.

\section{Results}

\begin{center}
\begin{tabular}{p{5.4cm}p{7.2cm}}
\toprule
\textbf{check} & \textbf{outcome} \\
\midrule
\texttt{\_patch\_grid(768$\times$768)}   & $(24,24)$ --- frozen grid reproduced \\
\texttt{\_patch\_grid(1024$\times$768)}  & $(32,24)$ --- adapts to input resolution \\
CI literal guard                         & passes (no forbidden literal in code) \\
full test suite                          & 110 passed \\
\bottomrule
\end{tabular}
\captionof{table}{Phase 1.7 verification. The patch grid reproduces the pilot's grid for the pilot's
input and adapts for a different one; the CI guard is active and green.}
\label{tab:p17}
\end{center}

The audit's effect is preventative rather than a headline number: the attribution heatmap now cannot
misalign if a scale-run model uses a different input resolution, the two remaining mislabel risks are
closed, and any future reintroduction of a dataset count or patch dimension as a code literal fails the
suite immediately --- the same class of bug Chapter~13 found by accident is now caught by construction.

\section{Standard vs. Non-Standard Choices}

\textbf{Standard.} Reading counts from data instead of hardcoding them, and asserting invariants at the
point they are relied upon, are ordinary defensive-programming practice.

\textbf{Non-standard, and the contribution.} Two choices. The patch grid is not merely moved to a
constant but \emph{derived from the actual input and guarded by an assertion}, so the code is correct by
construction across input sizes rather than correct by coincidence at one size --- directly answering
the plan's instruction to recompute the arithmetic rather than hardcode it. And the CI guard is
deliberately tokenizer-based so it can distinguish a magic literal in \emph{code} from the same number
appearing in a comment or docstring; a naive grep would either miss real regressions or drown in false
positives on the many places the report's prose legitimately cites 163 or 576. Encoding ``a literal
here is a bug, the same number in a sentence is documentation'' into the check is what makes it
maintainable enough to actually keep passing.

\section{Implications}

With the patch grid derived and asserted, Phase~4's attribution work and Phase~5's additional models can
run at whatever input resolution they require without a silent heatmap-misalignment risk. The CI guard
makes the ``every number traces to data or config'' principle enforceable rather than aspirational,
which matters most exactly when the codebase grows during the scale run and manual vigilance is least
reliable. This closes Phase~1: the six foundational fixes and this audit together leave a pipeline whose
metrics test the right region, measure model faithfulness rather than proxy plumbing where the fixes
reach, and reproduce the frozen pilot byte-for-byte under their default flags.
